% CAL-14: materiał kanoniczny w ME-23.
% Wszystkie 37 etykiet oryginalnego rozdziału zachowano w ME-23.
% Do wycofania książki pozostawiamy jedynie odsyłacze.
\providecommand{\atlasref}[2]{#2}
\section{Całka na grupie topologicznej: miara Haara}

Pełny wykład miary Haara, wraz z dowodami i przykładami,
znajduje się w \atlasref{ME-23::ch:haar-23}{rozdziale 23 teorii miary}.
Nie powtarzamy tu definicji ani konstrukcji.

\begin{itemize}
 \item \atlasref{ME-23::def:haar-23}{Lewe i prawe miary Haara};
 \item \atlasref{ME-23::thm:haar-main-23}{Istnienie i jednoznaczność bez założeń przeliczalności};
 \item \atlasref{ME-23::cor:haar-finiteness-23}{Skończoność i $\sigma$-skończoność};
 \item \atlasref{ME-23::thm:haar-modular-23}{Funkcja modularna};
 \item \atlasref{ME-23::cor:haar-integral-inversion-23}{Wzory całkowe i inwersja};
 \item \atlasref{ME-23::it:lem:haar-lie-volume}{Miara Haara na grupach Liego};
 \item \atlasref{ME-23::it:thm:haar-l1-convolution}{Algebra splotu $L^1(G)$};
 \item \atlasref{ME-23::it:ex:haar-ell2}{Znaczenie lokalnej zwartości};
 \item \atlasref{ME-23::tw:haar-automorfizm}{Moduł automorfizmu}.
\end{itemize}
Przykłady grup zwartych, dyskretnych, torusa,
grupy afinicznej i całek względem miary Haara
są omówione w tym samym rozdziale.
